
This appendix proves the two results the body argues from: the exchange
rate of Theorem \ref{thm:exchange}, whose exact pathwise form is stated as
Theorem \ref{thm:t2exact}, and the two-sided structure-proofness result of
Theorem \ref{thm:reach}, stated in full as Theorem \ref{thm:t9full}. Both
proofs are
reproduced from the derivation repository without change; the check
identifiers V2.1 and V9.1 to V9.4 refer to its verification suite. Proofs of
the remaining statements in Appendix \ref{app:statements} are in that
repository and are not reproduced here.

\subsection*{The cost-equivalence lemma}

\begin{proof}[Proof of Lemma L1 (cost equivalence)]
Taylor's theorem with Lagrange remainder gives, for each $t$,
\[
\Phi(\hstar) - \Phi(h_t)
 = -\delta_0\,\dev_t + \tfrac{1}{2}\gpo\,\dev_t^2 + r_t,
 \qquad |r_t| \le \tfrac{L_3}{6}\,|\dev_t|^3,
\]
with $L_3 = \sup |\Phi'''|$ on the operating interval. Summing over
$t \le T$ and taking expectations, it remains to show
$\E[\sum_t \dev_t] = 0$ under A2-sym. The recursion gives
$\E[\dev_{t+1} \mid \mathcal{F}_t] = -m\,\dev_t + \E[u_t \mid \mathcal{F}_t]
= -m\,\dev_t$, so $\E[\dev_{t+1}] = -m\,\E[\dev_t]$ by the tower property,
and $\E[\dev_t] = (-m)^t\,\E[\dev_0] = 0$ when the loop starts at the
operating point or in the stationary regime. The quadratic term contributes
$\tfrac{1}{2}\gpo\,\E[\sum_t \dev_t^2]$, which is the claim.

For the variance decomposition, write the realized cost (net of the cubic
remainder) as $C = -\delta_0 X + \tfrac{1}{2}\gpo Y$ with
$X = \sum_t \dev_t$, $Y = \sum_t \dev_t^2$. Then
$\Var(C) = \delta_0^2 \Var(X) + \tfrac{\gpo^2}{4}\Var(Y)
- \delta_0\gpo\,\Cov(X, Y)$. Under Gaussian innovations the vector
$(\dev_1,\dots,\dev_T)$ is jointly Gaussian with mean zero, and every third
moment of a zero-mean Gaussian vector vanishes; since
$\Cov(X,Y) = \sum_{t,s}\E[\dev_t\,\dev_s^2]$ is a sum of third moments, it is
zero. (For non-Gaussian symmetric innovations the same conclusion follows
from the sign symmetry $\dev \mapsto -\dev$ of the stationary law.)
Verified: V0.1--V0.3, including the necessity check that an asymmetric rule
breaks the expectation identity by exactly $-\delta_0\,\E[\sum_t \dev_t]$.
\end{proof}

\subsection*{The exchange rate}

\begin{proof}[Proof of Theorem T2 (pathwise exchange rate)]
Under A3, conditionally on the design $\{\dev_t\}$ the observations are
independent with means linear in $\eps$ and variance $\sigma^2$, so the
least-squares slope has conditional variance
$\Var(\hat\eps \mid \mathrm{design}) = \sigma^2/\Sxx$ (Gauss--Markov).
Multiplying by the quadratic cost $\tfrac{1}{2}\gpo \sum_t \dev_t^2$ and
using $\sum_t \dev_t^2 = \Sxx + T\bar\dev^2$ gives the displayed identity
exactly, on every realization. In the stationary regime
$\bar\dev = O_P(T^{-1/2})$ by the stationary long-run variance $\Var(\sum_{t\le T}\dev_t) = T v (1-m)/(1+m)(1+o(1))$, while
$\Sxx/T \to v$ a.s., so $T\bar\dev^2/\Sxx = O_P(1/T)$. Verified: V2.1
(machine precision).
\end{proof}


\subsection*{Structure-proofness, and the exactness of the reach}

\begin{proof}[Proof of Theorem T9 (structure-proofness within reach)]
Write $x = h - \hstar$, $d(h) = x$, and
$V = \mathrm{span}\{1,\, e^{-ch},\, h e^{-ch}\}$, the span of the
components of $s$ (a bijective correspondence $u \leftrightarrow
\phi_u = u^\top s$ between $\mathbb{R}^3$ and $V$, since $C_1 \neq 0$).
Differentiating $s$ gives
$s'(h) = (0,\, -c e^{-ch},\, C_1 e^{-ch}(ch - 1))^\top$, so
$g = -s'(\hstar)$ componentwise, which is the stated estimand gradient.

\emph{Step 1: one-dimensional reduction.} For $M(\mu) \succ 0$ the
generalized Rayleigh identity gives
\[
g^\top M(\mu)^{-1} g
 = \sup_{u \neq 0} \frac{(u^\top g)^2}{u^\top M(\mu)\, u}
 = \sup_{\phi \in V} \frac{\phi'(\hstar)^2}{\int \phi^2\, d\mu}
 = \Big(\inf\big\{ \textstyle\int \phi^2\, d\mu :
   \phi \in V,\ \phi'(\hstar) = 1 \big\}\Big)^{-1},
\]
using $u^\top M u = \int \phi_u^2\, d\mu$ and $u^\top g =
-\phi_u'(\hstar)$ (the sign is lost in the square). Hence
\[
R(\mu) \;=\; \frac{\int d^2\, d\mu}
{\inf\{\int \phi^2 d\mu : \phi \in V,\ \phi'(\hstar) = 1\}}.
\]

\emph{Step 2: the candidate and its pointwise domination.} Let
\[
\phi_0(h) \;=\; \frac{2}{c} - e^{-cx}\Big(\frac{2}{c} + x\Big)
 \;=\; \frac{2}{c}\cdot 1
 \;-\; e^{c\hstar}\Big(\frac{2}{c} - \hstar\Big)\, e^{-ch}
 \;-\; e^{c\hstar}\, \big(h e^{-ch}\big) \;\in\; V,
\]
the unique element of $V$ with 2-jet $(0, 1, 0)$ at $\hstar$ (the
2-jet map is invertible on $V$: in the basis $\{1, e^{-ch}, h e^{-ch}\}$
its determinant is $c^2 e^{-2c\hstar} \neq 0$); the expansion is
$\phi_0(x) = x - (c^2/6)\,x^3 + O(x^4)$. Writing
$f(x) = \tfrac{2}{c} - e^{-cx}(\tfrac{2}{c} + x)$ we claim
\[
|f(x)| \le |x| \quad \text{for all } x \ge -2/c,
\qquad \text{with equality iff } x \in \{0, -2/c\},
\]
and $|f(x)| > |x|$ strictly for every $x < -2/c$. Indeed:
(a)~$v(x) = f(x) - x$ has $v(0) = 0$ and
$v'(x) = e^{-cx}(1 + cx) - 1 < 0$ for $x \neq 0$ (since $1 + t < e^t$
for $t = cx \neq 0$); $v' \le 0$ with equality only at the isolated
point $0$ makes $v$ strictly decreasing on $\mathbb{R}$, so $f > x$ on
$x < 0$ and $f < x$ on $x > 0$.
(b)~$f(0) = 0$ and $f'(x) = e^{-cx}(1 + cx) > 0$ for $x \ge 0$, so
$f \ge 0$ on $[0, \infty)$ and, with (a), $0 \le f \le x$ there.
(c)~$q(x) = f(x) + x$ has $q(-2/c) = 0 = q(0)$ and
$q'(x) = e^{-cx}(1 + cx) + 1$, where $r(x) = e^{-cx}(1 + cx)$ is
strictly increasing on $x < 0$ ($r'(x) = -c^2 x\, e^{-cx} > 0$) from
$r(-2/c) = -e^2 < -1$ to $r(0) = 1$; so $q'$ changes sign exactly once
on $(-2/c, 0)$, from negative to positive, and $q \le 0$ there with
equality only at the endpoints; with (a) ($f \ge x$, equality only at
$0$) this gives $x \le f \le -x$ on $[-2/c, 0]$.
(d)~For $x < -2/c$, $r(x) < -e^2$ gives $q' < 0$, so moving left from
$-2/c$ the function $q$ increases above $0$: $f(x) > -x = |x| > 0$,
and in fact $f(x) = 2/c + e^{c|x|}(|x| - 2/c)$ grows exponentially.

\emph{Step 3: the bound.} If $\mathrm{supp}(\mu) \subset
[\hstar - 2/c, \infty)$, Step 2 gives $\int \phi_0^2\, d\mu \le \int
d^2\, d\mu$, so by Step 1, $R(\mu) \ge \int d^2 d\mu / \int \phi_0^2
d\mu \ge 1$. Strictness: equality in Step 2 holds only at
$x \in \{0, -2/c\}$, so any $\mu$ placing positive mass off
$\{\hstar, \hstar - 2/c\}$ (in particular any with three distinct
support points) has $\int \phi_0^2 d\mu < \int d^2 d\mu$ and
$R(\mu) > 1$. If $M(\mu)$ is singular, two cases:
$g \notin \mathrm{range}\,M(\mu)$ makes $\eps(\hstar)$ non-estimable
(infinite bound); $g \in \mathrm{range}\,M(\mu)$, which does occur for
special two-point designs (a one-parameter family of in-reach pairs;
their finite pseudo-inverse products were checked to obey the bound,
e.g.\ $R = 1.063$ at the pair $\{1.627, 2.127\}$), keeps
the sup-representation valid over $\{u : u^\top M u > 0\}$, and
$\int \phi_0^2 d\mu = 0$ would put $u_0 \in \mathrm{null}\, M$ with
$u_0^\top g = -1 \neq 0$, contradicting $g \in \mathrm{range}\,M$; so
$\int \phi_0^2\, d\mu > 0$ and the bound goes through verbatim.
Throughout, the Cram\'er--Rao functional bounds unbiased estimation,
and all regular estimators in the local-asymptotic-minimax sense.

\emph{Step 4: the infimum is 1.} The Cram\'er--Rao functional is
invariant under $C^1$ reparametrization, and the 2-jet map
$\theta \mapsto (\tau_\theta(\hstar),\, \eps(\hstar),\,
\tau_\theta''(\hstar))$ is a local diffeomorphism: its Jacobian
$[\,s(\hstar)\ {-s'(\hstar)}\ s''(\hstar)\,]$ has determinant
$C_1 c^2 e^{-2c\hstar} \neq 0$ (direct computation; V9.4). In jet
coordinates $\eta$ the model reads $\tau = \eta_1 - \eta_2 x +
\tfrac{1}{2}\eta_3 x^2 + O(x^3)$ with sensitivity $\tilde s(x) =
(1, -x, x^2/2)^\top + O(x^3)$, and the estimand is $\eta_2$. For the
symmetric three-point design $\mu_\delta$ uniform on $\{\hstar -
\delta, \hstar, \hstar + \delta\}$, scale by $D = \mathrm{diag}(1,
\delta, \delta^2)$: $M_\eta(\mu_\delta) = D\,\hat H_\delta\, D$ with
$\hat H_\delta \to \hat H_0$ invertible (the exact quadratic-model
moment matrix), so $(M_\eta^{-1})_{22} = \delta^{-2}(\hat
H_\delta^{-1})_{22}$ and $R(\mu_\delta) \to R_{\mathrm{quad}}$. In
the exact quadratic model, for \emph{any} symmetric design (odd
moments $m_1 = m_3 = 0$), the $(2,2)$ cofactor gives
$(M_\eta^{-1})_{22} = 1/m_2$ and $R_{\mathrm{quad}} = m_2 \cdot
(1/m_2) = 1$. Hence $R(\mu_\delta) \to 1$ (measured rate
$R(\mu_\delta) - 1 = O(\delta^2)$; V9.2), the infimum is 1, and by
Step 3 it is not attained.

\emph{Step 5: the reach is exact.} Below $\hstar - 2/c$ the
domination of Step 2 fails strictly, and the failure is realized by a
design, not only by the candidate: the frozen four-point witness
(support $\{1.8726, 1.8665, 1.3073, 0.0500\}$, weights
$\{0.003329, 0.955702, 0.040697, 0.000272\}$ after normalization, in
the register market $c = 1.5$, frozen anchor $\hstar = 1.8461$) has
$R = 0.851650\ldots$, verified at 50-digit precision (mpmath) and
reproduced in float64 by \texttt{run\_open1.py} (V9.2). All its
support except the probe at $0.05$ lies within reach: one
vanishing-weight below-reach probe breaks the floor. Symmetric
pair-plus-far-probe families do \emph{not} break it (their ratio
approaches 1 from above as the probe weight vanishes); the violation
requires the asymmetric cluster, which is why coarse searches missed
it. Verified: V9.1--V9.4.
\end{proof}

